% ------------------------------------------------------------------
%  CV — Manuel Barra Lazo (ES)
%  Compilar con pdflatex (Overleaf o TeX Live):  pdflatex Manuel_Barra_Lazo_CV_ES.tex
%  Una columna, texto seleccionable y sin iconos: pensado para ATS.
% ------------------------------------------------------------------
\documentclass[10pt,a4paper]{article}

\usepackage[T1]{fontenc}
\usepackage[utf8]{inputenc}
\usepackage{charter}
\usepackage[a4paper,top=1.1cm,bottom=0.9cm,left=1.45cm,right=1.45cm]{geometry}
\usepackage[protrusion=true,expansion=true]{microtype}
\usepackage{xcolor}
\usepackage{enumitem}
\usepackage{titlesec}
\usepackage{tabularx}
\usepackage[hidelinks]{hyperref}

\input{glyphtounicode}
\pdfgentounicode=1

\definecolor{ink}{HTML}{1A1A1A}
\definecolor{accent}{HTML}{0B5C7A}
\definecolor{muted}{HTML}{5A5A5A}

\hypersetup{
  pdftitle={Manuel Barra Lazo - CV},
  pdfauthor={Manuel Barra Lazo},
  pdfsubject={Desarrollador Full Stack},
  pdfkeywords={Full Stack, Frontend, React, Next.js, TypeScript, Node.js, Barcelona}
}

\pagestyle{empty}
\setlength{\parindent}{0pt}
\raggedright
\hyphenpenalty=10000
\exhyphenpenalty=10000
\color{ink}

\titleformat{\section}{\small\bfseries\color{accent}}{}{0pt}{\MakeUppercase}[\vspace{1pt}\titlerule]
\titlespacing*{\section}{0pt}{6pt}{3pt}

\setlist[itemize]{leftmargin=1.1em,label={\color{accent}\textbullet},topsep=2pt,itemsep=0.6pt,parsep=0pt}

% \role{Empresa}{Puesto}{Fechas}{Ubicación}
\newcommand{\role}[4]{%
  \begin{tabularx}{\linewidth}{@{}X r@{}}
    \textbf{#1} \,·\, #2 & \small\color{muted}#4 \,·\, #3
  \end{tabularx}\par\vspace{0pt}}
\newcommand{\context}[1]{{\small #1}\par}
\newcommand{\stack}[1]{\vspace{1pt}{\footnotesize\color{muted}\textbf{Stack:} #1}\par\vspace{4pt}}

\begin{document}

% ---------------- Cabecera ----------------
{\LARGE\bfseries Manuel Barra Lazo}\par\vspace{2pt}
{\large\color{accent} Desarrollador Full Stack \,·\, Next.js, React y TypeScript}\par\vspace{3pt}
{\footnotesize Barcelona \,·\, +34 634 008 260 \,·\,
\href{mailto:manuelbarralazo@gmail.com}{manuelbarralazo@gmail.com} \,·\,
\href{https://www.malb.dev}{malb.dev} \,·\,
\href{https://www.linkedin.com/in/manuelbarra}{linkedin.com/in/manuelbarra} \,·\,
\href{https://github.com/ManuelBarra}{github.com/ManuelBarra}}

% ---------------- Perfil ----------------
\section{Perfil}
Llevo más de cinco años haciendo producto web con React y Next.js. Empecé en consultoría, con proyectos para CaixaBank, Roca y Repsol, y después pasé a Genomcore, una healthtech, donde llevé sus dos webs de WordPress a Next.js y metí la IA en el día a día del equipo de marketing. Ahora trabajo como desarrollador full stack en SADATA y, en paralelo, estudio DAM para sumar Java, Kotlin y Android.

% ---------------- Experiencia ----------------
\section{Experiencia}

\role{SADATA}{Desarrollador Multiplataforma (Full Stack)}{jun 2026 -- actualidad}{Remoto}
\context{Consultora de BI, IA y software a medida, con equipo en Barcelona y Santiago de Chile y clientes en 5 países.}
\begin{itemize}
  \item Estoy terminando \textbf{Hip Hop 4 Life}, nuestra plataforma para aprender cultura hip hop: 134 lecciones en 6 disciplinas, progresión por niveles y suscripciones con Stripe. La he construido con Next.js 16, TypeScript, Prisma y PostgreSQL.
  \item Le añadí IA con la API de Claude: analiza la métrica y la rima de cada verso y da feedback a los alumnos, con límites de uso para evitar abusos.
  \item Para \textbf{BDI}, una plataforma de email marketing de Reino Unido, escribí la especificación de producto de «Source Contacts» y desarrollé su frontend en React (jun--ago 2026).
\end{itemize}
\stack{TypeScript, Next.js, React, Node.js, Prisma, PostgreSQL, Anthropic SDK, Stripe, Jest, Vercel}

\role{Genomcore}{Frontend Developer · IA aplicada}{ene 2025 -- ago 2026}{Barcelona}
\context{Healthtech con más de 10 años trabajando con datos clínicos y genómicos para hospitales y servicios de salud.}
\begin{itemize}
  \item Migré \textbf{genomcore.com} y \textbf{madeofgenes.ai} de WordPress a Next.js. Las dos webs sacan \textbf{100 en las cuatro categorías de Lighthouse}.
  \item Rediseñé ambas junto a un diseñador, incluida la web del 10.º aniversario de la empresa.
  \item Creé un \textbf{catálogo de skills de IA} para el equipo de marketing (7 personas) que les ahorra unas \textbf{10 horas a la semana}.
  \item Usé Claude Code para acelerar la migración del código antiguo.
\end{itemize}
\stack{Next.js, React, TypeScript, WordPress, Markdown, Claude Code, Gemini, Git}

\role{DigitalMakers (Vass)}{Frontend Developer}{abr 2022 -- jun 2024}{Barcelona}
\begin{itemize}
  \item Me encargué de todo el frontend de \textbf{AgroBank Hub} (CaixaBank): agenda de eventos, casos de innovación, tienda e informes, todo conectado a APIs REST.
  \item Desarrollé y mantuve los portales de \textbf{CaixaBank} (Talks, Banca Privada, Banca Premier y Wealth Management) y sus versiones en 6 países.
  \item Trabajé en el frontend de \textbf{Roca.es} y \textbf{Guía Repsol}, desde la maquetación hasta las incidencias en producción.
\end{itemize}
\stack{JavaScript, Liferay, AEM, jQuery, SASS, Bootstrap, REST APIs, Figma, Jira, Scrum}

{\small\color{muted}\itshape Pausa profesional para viajar, jul -- dic 2024.}\par\vspace{4pt}

\role{inMotion Real Estate Media}{Frontend Developer (en inglés)}{nov 2021 -- abr 2022}{Barcelona}
\begin{itemize}
  \item Construí entera la web de \textbf{The Abbey Company}, con un buscador de propiedades en Vue.js (filtros, mapa, cuadrícula y lista), vacantes y estadísticas.
  \item Hice temas de WordPress a medida para The Multifamily Firm y H\&R Retail, dos inmobiliarias de EE.\,UU.
\end{itemize}
\vspace{4pt}

\role{Clos del Portal · Vins Nus · Alfredo Arribas}{Desarrollador web freelance}{sep 2020 -- actualidad}{Remoto, en paralelo}
\begin{itemize}
  \item Desde 2020 me ocupo de las webs y la tienda online de un grupo de bodegas del Priorat (hosting, seguridad y SEO), y ahora las estoy pasando de WordPress a Next.js.
\end{itemize}

% ---------------- Formación ----------------
\section{Formación}
\begin{tabularx}{\linewidth}{@{}X r@{}}
\textbf{Técnico Superior en Desarrollo de Aplicaciones Multiplataforma (DAM)} \,·\, Jobie {\small(en curso)} & \small\color{muted}2026 -- 2028 \\[2pt]
\textbf{Bootcamp Frontend Developer} \,·\, Skylab Coders Academy (ISDI Coders) & \small\color{muted}2021 \\[2pt]
\textbf{Desarrollo y maquetación web} \,·\, Espai Escola Professional d'Aplicacions Informàtiques & \small\color{muted}2019 -- 2021 \\
\end{tabularx}

% ---------------- Habilidades ----------------
\section{Habilidades}
\begin{tabularx}{\linewidth}{@{}l X@{}}
\textbf{Frontend} & TypeScript, JavaScript, React, Next.js, Vue.js, Tailwind CSS, SASS, GSAP, HTML/CSS \\
\textbf{Backend y CMS} & Node.js, Express, Prisma, PostgreSQL, MongoDB, PHP, WordPress, Liferay, AEM \\
\textbf{IA} & Claude API (Anthropic SDK), Claude Code, Gemini \\
\textbf{Workflow} & Git/GitHub, Jest y Testing Library, Vercel, Lighthouse y Core Web Vitals, SEO, Figma, Jira, Scrum \\
\textbf{Idiomas} & Español (nativo) · Inglés (avanzado; he trabajado en inglés con clientes de EE.\,UU. y Reino Unido) \\
\end{tabularx}

\end{document}
